\documentclass[11pt]{article}
\usepackage[margin=1in]{geometry}
\usepackage{amsmath, amssymb, bm}
\usepackage{parskip}

\begin{document}

\section*{EKF Algorithm and Derivation}

An EKF works by keeping a best guess of the robot's current state, $\bm{x}$, and of how
uncertain that guess is, $P$. Each iteration predicts where the robot moved, then corrects
that prediction using sensor data. The prediction and the measurements are weighted by how
reliable each one is.

Here, the state is $\bm{x} = [x,\ y,\ \theta,\ \omega,\ v,\ \dot v]^T$ and the
measurements are odometry ($v,\omega$) and IMU ($a_x,a_y$) readings,
$\bm{z} = [v,\ \omega,\ a_x,\ a_y]^T$. $P$ is the $6\times6$ state covariance. Its diagonal
holds the variance of each state, and its off-diagonal terms describe how errors in
different states are correlated. Each iteration ($dt = 0.1$\,s) has two steps.

\subsection*{Prediction}
\begin{equation*}
\bm{x} = f(\bm{x}), \qquad P = APA^T + Q
\end{equation*}
The motion model $f$ moves the state forward one time step. $A$, its Jacobian, carries
the existing uncertainty forward with the motion, and the process noise $Q$ adds
uncertainty to account for model error. Prediction therefore always increases $P$.

\textbf{Derivation.} We use unicycle kinematics, assuming $\omega$ and $\dot v$ are
constant over $dt$ ($\theta$ is wrapped to $[-\pi,\pi]$). Differentiating each row of
$f$ with respect to each state gives $A = \partial f/\partial\bm{x}$:
\begin{equation*}
f(\bm{x}) =
\begin{bmatrix}
x + v\cos\theta\,dt \\ y + v\sin\theta\,dt \\ \theta + \omega\,dt \\ \omega \\ v + \dot v\,dt \\ \dot v
\end{bmatrix},
\qquad
A =
\begin{bmatrix}
1 & 0 & -v\sin\theta\,dt & 0 & \cos\theta\,dt & 0\\
0 & 1 & v\cos\theta\,dt & 0 & \sin\theta\,dt & 0\\
0 & 0 & 1 & dt & 0 & 0\\
0 & 0 & 0 & 1 & 0 & 0\\
0 & 0 & 0 & 0 & 1 & dt\\
0 & 0 & 0 & 0 & 0 & 1
\end{bmatrix}
\end{equation*}
The off-diagonal terms in the first two rows link heading and speed errors to position
errors, which builds the correlations in $P$ that the update relies on.

\subsection*{Update}
\begin{align*}
\bar{\bm{y}} &= \bm{z} - h(\bm{x}), & S &= CPC^T + R, & K &= PC^TS^{-1}, \\
\bm{x} &= \bm{x} + K\bar{\bm{y}}, & P &= (I - KC)P
\end{align*}
The measurement model $h$ gives the readings expected at the predicted state, so the
innovation $\bar{\bm{y}}$ is the difference between what the sensors read and what was
expected. $S$ is the expected size of that difference: the state uncertainty, mapped into
measurement space by the Jacobian $C$, plus the measurement noise $R$. The Kalman gain
$K$ is the ratio of state uncertainty to total uncertainty. If $R$ is large relative to
$P$, $K$ is small and the filter trusts the model. If $P$ dominates, the filter follows
the sensors. The state is corrected by $K\bar{\bm{y}}$, and $P$ shrinks.

\textbf{Derivation.} Odometry measures $v$ and $\omega$ directly. The IMU's forward
acceleration is $\dot v$, and its lateral acceleration is the centripetal term
$a_y = v^2/r = v\omega$ (using $r = v/\omega$). This gives $h$, and $C = \partial h/\partial\bm{x}$:
\begin{equation*}
h(\bm{x}) =
\begin{bmatrix} v \\ \omega \\ \dot v \\ v\omega \end{bmatrix},
\qquad
C =
\begin{bmatrix}
0 & 0 & 0 & 0 & 1 & 0\\
0 & 0 & 0 & 1 & 0 & 0\\
0 & 0 & 0 & 0 & 0 & 1\\
0 & 0 & 0 & v & \omega & 0
\end{bmatrix}
\end{equation*}
The columns of $C$ for $x$, $y$ and $\theta$ are zero, because no sensor measures them.
They are corrected only through their correlations in $P$, so their uncertainty grows
over time.

\subsection*{Linearization and tuning}
$f$ (through $\cos\theta$, $\sin\theta$) and $h$ (through $v\omega$) are nonlinear, so the
EKF linearizes them about the current estimate at every step. The exact functions
propagate the state, and the Jacobians propagate the covariance. $A$ is evaluated at the
previous state and $C$ at the predicted state.

The filter starts at $\bm{x}_0 = \bm{0}$ with $P_0 = I_6$, and uses
$Q = q I_6$ and $R = r I_4$. Using scaled identity matrices assumes independent noise of
equal size in every state and sensor. The ratio $Q/R$ sets the tradeoff: a larger ratio
tracks the measurements closely but is noisier, and a smaller ratio is smoother but lags.
The final tuning was $q = 0.6$, $r = 0.15$.

\end{document}
